% Français 9115, batch 16 — Gallica views 301–320.
% Pass: Opus 5.5 (claude-opus-5-5), 2026-09-22 — first pass, unchecked against the views by a human.
%
% What the twenty views hold. Ten views carry text, and they are nine
% leaves: every leaf is written on the recto only, and the views between
% them (303, 305, 307, 309, 311, 313, 315, 317, 319) are blank versos
% showing nothing but the recto's writing through the paper, reversed.
% View 302 is not a new leaf but a second exposure of the leaf of view
% 301 (same text, same ink numbers, framing slightly different); it is not
% transcribed. There is no microfilm target. Three pieces follow one
% another, and the batch holds the end of one piece and the beginning of
% another.
%
%   v. 301–308  The end of her French memoir examining Euler's § 47 (« N.o
%               47 »), which began at view 218: her pagination at the head
%               runs on, 39 to 42, and her paragraph numbers reach 24.
%               View 301 continues the sentence that stops view 299
%               (« …par cause que les deux / premières de ces équations sont
%               semblables… »). § 24 (v. 301): the rod free at both ends
%               with an odd number of nodes may be regarded as two rods half
%               as long, free at one end and supported at the other, and
%               likewise for fixed ends — a remark independent of Euler's
%               four equations, got from y = 0 at x = 1/2 alone; so Euler's
%               six cases reduce to four « reellement independans », cases
%               I, III, IV and VI (his §§ 15, 30, 35, 43). V. 304–306: the
%               check on the coefficients — case II against case I agrees
%               once 1 is put for 1/2; case V against case VI seems not to,
%               until one observes that Euler puts x = 0 at the supported
%               end, and substitutes 1−x for x; the ordinate's second
%               derivative vanishing at the middle confirms that the middle
%               behaves as a supported end (Euler's N.o 13). V. 308: hence
%               the coefficients computed for the whole rod serve for the
%               rods it may be divided into; the free-free rod with no node
%               at the middle cannot be so split. The text stops short on
%               « extremités », the foot of the leaf blank. Fair hand,
%               lightly corrected.
%   v. 310–316  Still paginated 43, 44, 45 at the head, but a draft, and a
%               new subject: from the transverse vibrations of the lamina
%               to square elastic plates vibrating in several directions,
%               after Chladni (« traité d'Acoustique par Chladni §§ 110 »,
%               his figures 63, 64 a, 75, 82). A plate whose nodal figure
%               has no line cutting it diametrically in two equal parts can
%               be composed of four (nine, sixteen) plates bearing the same
%               figure; if a plate aamm is divided into mm squares aa and a
%               stylet makes a figure in one of them, the figure repeats in
%               the others, the nodal curves cutting the horizontal plane at
%               α, α+a, α+2a … and α', α'+a, α'+2a …; the analogy with
%               tensed vibrating surfaces, citing « § 26 du mémoire de M.r
%               Biot » in « le tome 4 des memoires de l'institut », where
%               α = α' = 0. V. 312 and 314 are heavily corrected, with long
%               continuations in the left margin and passages cancelled by
%               hatching. V. 314: the complete theory of the square plate
%               would have fifteen cases, according as the sides are
%               supported, free or fixed; irregular figures suggest that the
%               point of support hinders the motion more than the other
%               nodal points. V. 316, unpaginated at the head, is another
%               draft of the opening of v. 314, going on to the « immensité
%               de cas » made « effrayante » by the different degrees of
%               pressure Chladni reports: a weak pressure differs little
%               from freedom, and fixity is the limit of an infinite scale.
%   v. 318      A title leaf, « Memoire presenté et examiné / sur la lame »,
%               and nothing else.
%   v. 320      A new piece in fair copy, its own page (1):
%               « Remarques sur le mémoire d'Euler: Investigatio motuum
%               quibus laminæ e virgæ elasticæ contremiscunt. in act acad
%               Petrop ann. 1779 P 1 p. 103 et seq », a summary in the third
%               person of Euler's § 47 and following (the equation for ω,
%               the case x = 1/2, the two solutions of §§ 49–51, cases IV
%               and V). The text stops at « au cas V. », two thirds down the
%               leaf; whether it goes on after view 321 is not in this batch.
%
% The numbers on the leaves, and why no \folio{} is written. (1) At the
% top right of every written leaf, in the larger, rounder hand and broader
% pen of the running ink series: 147 (v. 301, repeated on its second
% exposure v. 302), 148 (304), 149 (306), 150 (308), 151 (310), 152 (312),
% 153 (314), 154 (316), 155 (318), 156 (320) — one per leaf, without a gap,
% continuing 146 at view 299. It stands where the Bibliothèque's foliation
% stands but is penned, not pencilled, and following the house rule of this
% volume no \folio{} is written. (2) At the head, in the middle, underlined,
% in her own hand: 39 (301), 40 (304), 41 (306), 42 (308), 43 (310), 44
% (312), 45 (314) — her pagination of the memoir, continuing 38 at view 299;
% v. 316 and the title leaf v. 318 carry none; v. 320 begins a new series
% at « (1) ». (3) Her paragraph number 24, in the left margin of v. 301,
% continuing 23 at view 299; no other paragraph number in the batch. Her
% references to Euler's sections (§§ 13, 15, 23, 30, 35, 39, 43, 47, 49–51)
% and to Chladni's (§§ 110, figures 63–82) are text, not numbering. Every
% number above was read on its view; none is inferred. The blank versos
% carry no number of their own.
%
% Notation. Hers, as in batch 15: $x$ a fraction of the length, $\omega$
% the angle of the cases, $\alpha\beta\gamma\delta$ and $\alpha'\ldots$
% Euler's coefficients (her $\gamma$ written like a « v »); « sin », « cos »,
% « tang » set as \text{sin} etc.; her braces $\{\ \}$ kept; « $\ldots$ »
% for the dots she writes for an omitted factor in $y = \ldots$. In the
% plate pages, $a$ is the side of an elementary square and $m$ their number
% along a side.
%
% Spelling is hers and is not flagged: « coëfficiens », « equivallante »,
% « calculler », « plainement », « fesant », « apparans », « precedant »,
% « complette », « quarré » beside « carré », « analisées », « etres »,
% « leurs » for « leur », « faits » for « faite », « euler » without
% capital (v. 316), and the run-together small words — dela, endeux,
% demême, lamoitié, delaquelle, cequi, alaverité, degrandeurs, aucentre.
% Where the leaf and the calculation disagree the leaf stands, with a note
% (v. 301, the sign of tang ½ω; v. 316, « de dans » for « de n'en »).
%
% Her way of correcting, as settled by earlier passes: interlinear
% additions as \add{} beside the \struck{} they replace, in the order the
% sense suggests; long joining strokes not taken for deletions unless a
% separate heavier stroke runs through; passages cancelled by oblique
% hatching (v. 314) given as \struck{} and said so in a note. The long
% notes in the left margin of v. 312, 314 and 316 are hers and most of them
% complete the line they face; they are given as \marginal{} after the
% passage they belong to.
%
% What could not be read. \ill{} stands in twelve places: a short struck
% word after « lorsque » (v. 301); a heavily inked struck group after
% « cas VI », another after « sont » struck at « cas V », and a small sign
% before « $x = 0$ » (v. 304); a struck group after $\alpha+2a$, a line
% opening of struck terms, and the last word of a struck reference to Biot
% (v. 312), a struck word in the margin note after « demontrer » (v. 312);
% the struck word opening the second paragraph and the word opening the last
% struck line (v. 314); a struck sign after « de dans » and a struck word
% after « la fixité que » (v. 316). Readings offered with doubt are marked
% \uncertain{}: « $A = -e^{\omega}$ », « $x = 1$ » and « l'hypothese »
% (v. 304); « et » (v. 306); « $mm-1$ », « \&c. » and « § 26 » (v. 312);
% « examiné » (v. 318). Nothing was guessed to fill a gap.
%
% Nothing here was checked against any published edition of Euler, Chladni
% or Biot, or of Sophie Germain's papers.
\documentclass[11pt,a4paper]{article}
% The preamble shared by every transcription of the mathematicians' manuscripts in Gallica.
%
% It rests on one idea: **uncertainty compounds, it does not resolve.** A
% transcription that smooths over a doubtful word has destroyed the only thing
% it was made for — a reader must be able to tell, at every point, what was on
% the page and what was supplied. The apparatus macros below are therefore the
% critical apparatus, not decoration.
%
% The same commands are recognised by `scripts/render.mjs`, which produces the
% site's reading view, and by `scripts/tei.mjs`. A macro added here must be
% added there too, or rendering fails loudly — which is the wanted behaviour.
%
% Comments here are English, like the rest of the repository. The documents
% this preamble sets are French, and that is deliberate: see the skills.

\ProvidesPackage{germain}[2026/09/15 Transcriptions of mathematicians' manuscripts in Gallica]

\usepackage{amsmath,amssymb,amsthm}
\usepackage{mathrsfs}
% Kept from the parent preamble so the renderer's diagram support stays
% available; these manuscripts draw no commutative diagrams, and nothing here uses it.
\usepackage{tikz-cd}
\usepackage[english,french]{babel}
\usepackage{fontspec}

% --- Greek ------------------------------------------------------------------
% The text face stays fontspec's default, Latin Modern, which has no Greek:
% XeTeX drops every glyph it lacks with a line in the log and nothing on the
% page, and the Greek words the manuscripts quote vanished from the PDFs. So
% the Greek and Greek Extended blocks get a XeTeX character class of their own,
% and entering or leaving a run of it switches the family to CMU Serif — the
% same Computer Modern design, with polytonic Greek — and back. Only the
% family changes: a Greek word in \textit or \textbf keeps its shape and series.
% The transitions to and from every other class carry the tokens that class
% has towards an ordinary letter, so French spacing before « ; » or inside
% « » still holds after a Greek word. No grouping, so no brace or \endgroup
% can come out unbalanced; maths is left alone, since XeTeX inserts nothing
% there. Kept identical in `exercices.sty`.
\makeatletter
\newfontfamily\germainGreekfont{cmun}[
  Extension=.otf, NFSSFamily=germaingreek,
  UprightFont=*rm, ItalicFont=*ti, SlantedFont=*sl,
  BoldFont=*bx, BoldItalicFont=*bi, BoldSlantedFont=*bl]
\newXeTeXintercharclass\germain@greek
\def\germain@greekrange#1#2{%
  \@tempcnta=#1\relax
  \loop
    \XeTeXcharclass\@tempcnta=\germain@greek
    \ifnum\@tempcnta<#2\advance\@tempcnta\@ne\repeat}
\germain@greekrange{"0370}{"03FF}
\germain@greekrange{"1F00}{"1FFF}
\def\germain@greekprev{\rmdefault}
\def\germain@greekin{%
  \edef\germain@greekprev{\f@family}\fontfamily{germaingreek}\selectfont}
\def\germain@greekout{\fontfamily{\germain@greekprev}\selectfont}
\def\germain@greekjoin{%
  \XeTeXinterchartokenstate=\@ne
  \@tempcnta=\z@
  \loop
    \ifnum\@tempcnta=\germain@greek\else
      \XeTeXinterchartoks\germain@greek\@tempcnta=\expandafter{%
        \expandafter\germain@greekout\the\XeTeXinterchartoks\z@\@tempcnta}%
      \XeTeXinterchartoks\@tempcnta\germain@greek=\expandafter{%
        \the\XeTeXinterchartoks\@tempcnta\z@\germain@greekin}%
    \fi
    \ifnum\@tempcnta<4095\advance\@tempcnta\@ne\repeat}
% babel-french sets its own classes, and saves and restores the interchar
% state, each time a language is selected: join again after it does.
\addto\extrasfrench{\germain@greekjoin}
\addto\extrasenglish{\germain@greekjoin}
\AtBeginDocument{\germain@greekjoin}
\makeatother

\usepackage{geometry}
\geometry{a4paper,margin=25mm,marginparwidth=28mm}
\usepackage{marginnote}
\usepackage{xcolor}
\usepackage[normalem]{ulem}
\setlength{\emergencystretch}{3em}
\usepackage{hyperref}
\hypersetup{colorlinks=true,linkcolor=black,urlcolor=black,pdfborder={0 0 0}}

% --- Batch metadata ---------------------------------------------------------
% Taken as they stand by the rendering script, which shows them at the head of
% the reading view. They fix what the file claims to cover. The macro names are
% the parent project's, so that the scripts stay shared; \folder here means the
% volume's slug — `fr-9115` — and \pages counts Gallica views.

\newcommand{\folder}[1]{\gdef\grothFolder{#1}}
\newcommand{\batch}[1]{\gdef\grothBatch{#1}}
\newcommand{\pages}[2]{\gdef\grothFirst{#1}\gdef\grothLast{#2}}
\newcommand{\foldertitle}[1]{\gdef\grothFoldertitle{#1}}

% The holder's own shelfmark, printed as they write it: « Français 9115 ».
\newcommand{\shelfmark}[1]{\gdef\grothShelfmark{#1}}

% The dating, copied verbatim from the holder's catalogue — never inferred
% here. « XIXe siècle » is the BnF's; a pass that dates a leaf from its content
% says so in a \note{}, as its own inference.
\newcommand{\dating}[1]{\gdef\grothDating{#1}}

% The status notice, printed in the title block and repeated as a diagonal
% watermark on every page of the PDF. The manuscripts are in the public
% domain; what this line declares is not a right but a quality — an
% unverified first pass by a machine. The skills write the line; a file
% without it gets no watermark, so leaving it out is visible in the output.
\newcommand{\watermark}[1]{\gdef\grothWatermark{#1}}

\gdef\grothFolder{?}\gdef\grothBatch{}\gdef\grothFirst{?}\gdef\grothLast{?}%
\gdef\grothFoldertitle{}\gdef\grothShelfmark{}\gdef\grothDating{}\gdef\grothWatermark{}

% --- The résumé -------------------------------------------------------------
\newenvironment{resume}
  {\small\setlength{\parskip}{0.45em}}
  {\par\medskip\hrule\medskip}

% --- Keywords ---------------------------------------------------------------
% In English, closing the modernised reading's résumé: the modern vocabulary
% under which to search for what the leaves construct. The manifest extracts
% them to tag the volume — this line is the single source of the tags.
\newcommand{\keywords}[1]{\par\smallskip\noindent
  {\footnotesize\textsc{Keywords} --- \textit{#1}}\par}

% --- The critical apparatus -------------------------------------------------

% The Gallica view number, in the margin. This is the anchor that lets a
% disagreement over a formula be settled by looking at *one* image rather than
% a volume of 750. The site uses it to turn the facsimile as the transcription
% is scrolled. A view is one scanned image: usually one side of a leaf.
\newcommand{\page}[1]{%
  \marginnote{\footnotesize\color{gray!70}\textsf{v.\,#1}}%
  \label{page:#1}\ignorespaces}

% The BnF's pencilled foliation, where the leaf carries it and a pass can read
% it: « 198r ». This is what the literature cites — « ff. 198r–208v » — and
% the one line that turns a citation into a view. Recorded, never inferred:
% a folio a pass cannot read is not written.
\newcommand{\folio}[1]{%
  \marginnote{\footnotesize\color{gray!50}\textsf{f.\,#1}}\ignorespaces}

% The modernised reading's coarser counterpart to \page{}: which views a
% section is drawn from.
\newcommand{\pagerange}[2]{%
  \marginnote{\footnotesize\color{gray!70}\textsf{v.\,#1--#2}}%
  \ignorespaces}

% Illegible: never guessed.
\newcommand{\ill}{\textcolor{red!60!black}{[\,\ldots\,]}}

% A doubtful reading: the word is given, and flagged as uncertain.
\newcommand{\uncertain}[1]{\uline{#1}}

% An editorial addition — a missing letter, an implied word.
\newcommand{\add}[1]{\textcolor{blue!50!black}{[#1]}}

% Struck out by the writer. What was crossed out is part of the document.
\newcommand{\struck}[1]{\sout{#1}}

% The transcriber's note, on the state of the page or the mathematics.
\newcommand{\note}[1]{\par\smallskip{\small\color{gray!60!black}\textit{[#1]}}\par\smallskip}

% A marginal note of hers, distinct from the body of the text.
\newcommand{\marginal}[1]{\par\smallskip\noindent\hspace{1em}%
  {\small\color{gray!60!black}#1}\par\smallskip}

% --- Title ------------------------------------------------------------------
% The title block is French, like the documents themselves.

\AddToHook{shipout/background}{%
  \ifx\grothWatermark\empty\else
    \begin{tikzpicture}[remember picture, overlay]
      \node[rotate=52, scale=3.4, align=center, text=gray!18,
            font=\sffamily\bfseries]
        at (current page.center) {\grothWatermark};
    \end{tikzpicture}%
  \fi}

\AtBeginDocument{%
  \begin{center}
    {\large\bfseries Manuscrits de mathématiciens (Gallica) ---
      \ifx\grothShelfmark\empty\grothFolder\else\grothShelfmark\fi}\\[2pt]
    {\itshape\grothFoldertitle}\\[4pt]
    {\small vues \grothFirst--\grothLast
      \ifx\grothBatch\empty\else\ (lot \grothBatch)\fi}\\[2pt]
    \ifx\grothDating\empty\else
      {\small\color{gray!60!black}Datation du catalogue : \grothDating}\\[2pt]
    \fi
    \ifx\grothWatermark\empty\else
      {\small\bfseries\color{red!45!gray}\grothWatermark}\\[2pt]
    \fi
    {\footnotesize\color{gray!60!black}Transcription établie d'après les images IIIF de Gallica.
     Source gallica.bnf.fr / Bibliothèque nationale de France.\\
     Les lectures incertaines sont soulignées, les illisibles notées [\,\ldots\,],
     les ajouts entre crochets ; \textsf{v.} numérote les vues de Gallica,
     \textsf{f.} la foliotation au crayon de la BnF.}
  \end{center}
  \vspace{1em}}

\endinput


\folder{fr-9115}
\shelfmark{Français 9115}
\batch{16}
\pages{301}{320}
\dating{1801-1900}
\watermark{Lecture automatique\\non vérifiée}
\foldertitle{Papiers de Sophie GERMAIN. — Recueil de dissertations et problèmes mathématiques et physiques.}

\begin{document}

\note{Numérotation des feuillets. Chaque feuillet écrit porte en haut à
droite un nombre à l'encre, d'une plume large et d'une écriture ronde :
147 (vue 301, et de nouveau à la vue 302, seconde prise du même feuillet),
148 (304), 149 (306), 150 (308), 151 (310), 152 (312), 153 (314), 154
(316), 155 (318), 156 (320), un par feuillet et sans lacune. Cette série
occupe la place de la foliotation de la Bibliothèque, mais elle est tracée
à la plume et non au crayon ; aucun « f. » n'est donc porté. Au milieu de
la tête, soulignée, sa propre pagination : 39 (301), 40 (304), 41 (306),
42 (308), 43 (310), 44 (312), 45 (314) ; rien aux vues 316 et 318 ;
« (1) » à la vue 320, où commence une autre pièce. Tous ces nombres ont
été lus sur la vue ; aucun n'est déduit. Les vues 303, 305, 307, 309, 311,
313, 315, 317 et 319 sont les versos blancs des feuillets qui les
précèdent et ne sont pas transcrites ; la vue 302 reprend la vue 301.}

\page{301}

\struck{qu} premières de ces équations sont semblables entr'elles et se
reduisent l'une et l'autre a $y = 0$ et que les deux dernières
sont identiques, \struck{ou} lorsque \struck{\ill{}} les coëfficiens $\alpha, \beta, \gamma, \delta$ et $\alpha', \beta', \gamma', \delta'$
sont semblables entr'eux.

24 Deplus j'ai déja remarqué que si l'on pouvoit considérer \struck{tout}
la lame dont les deux extremités sont libres \add{lorsqu'elle a un nombre impair de nœuds,} comme partagée
endeux autres lames moitiés moins longues qui auroient chacune une
de leurs extremités libres et l'autre appuiée. et que demême aussi
la lame dont les deux extremités sont fixées lorsqu'elle a pareillement
un nombre impair de nœuds peut être censée partagée en deux
autres lames moitié moins longues qui auroient chacune une de leurs
extremités fixées et l'autre appuiée, Cette remarque est également
independante des quatre équations d'Euler et resulte tout
simplement de la supposition $y = 0$ pour $x = \frac{1}{2}$ car elle
donne pour l'un et l'autre cas $(1+e^{\omega})\,\text{sin}\,\frac{1}{2}\omega + (1-e^{\omega})\,\text{cos}\,\frac{1}{2}\omega = 0$
d'ou on tire $\text{tang}\,\frac{1}{2}\omega = \frac{1-e^{\omega}}{1+e^{\omega}}$ qui est la condition \struck{delaquelle} qui a été
deduite. dela considération de l'équation du mouvement de
la lame assujettie a l'action d'un stilet.

Cette maniere de considerer la lame comme partagée endeux autres
sous les conditions indiquées mène a conclure toute les circonstances
du mouvement d'une lame dont une extremité est libre et
l'autre appuiée de celles qui appartiennent au mouvement
dela lame dont les deux extrémités sont libres et demême
aussi cequi concerne le mouvement dela lame dont une extremité
est fixée et l'autre appuiée des valeurs relatives au mouvement
dela lame dont les deux extremités sont fixées. Desorte que
les six cas calculés par Euler peuvent etre reduits a 4
reellement \add{independans les uns des autres} savoir le cas I ou les deux extremités sont libres
§ 15 du memoire d'Euler le cas III ou une extremité est fixée et
l'autre libre § 30. le cas IV ou les deux extremités sont appuiées
§ 35 enfin le cas VI ou l'une et l'autre extremités sont fixées
§ 43
\note{En tête, au milieu, « 39 » souligné, de sa main ; en haut à droite
« 147 », à l'encre, de la série des lots précédents. Le paragraphe porte
le numéro 24, dans la marge de gauche. La première ligne poursuit la
phrase qui clôt la vue 299 (« par cause que les deux ») : « les deux /
premières de ces équations » ; le « qu » biffé qui ouvre la ligne est
d'une plume plus lourde. Un mot court biffé après « lorsque » n'est pas
lu. L'addition « lorsqu'elle a un nombre impair de nœuds, » est écrite
au-dessus de la ligne, d'une plume plus fine, avec un petit signe
d'insertion après « libres » ; « independans les uns des autres » de même
au-dessus de « savoir ». « moitiés » est suivi d'un petit trait en forme
de croix, sans renvoi visible. La valeur de $\text{tang}\,\frac{1}{2}\omega$
est écrite $\frac{1-e^{\omega}}{1+e^{\omega}}$, alors que l'équation qui
précède donne $-\frac{1-e^{\omega}}{1+e^{\omega}}$ ; la leçon du feuillet
est gardée. Après « deduite. » la ligne laisse un blanc d'environ un
tiers de sa longueur avant « dela considération » ; la phrase se lit sans
coupure. Les deux derniers renvois sont lus « § 35 » et « § 43 » : le
second chiffre de « 35 » est le même tracé à queue descendante que celui
de « § 15 » deux lignes plus haut, et ces numéros sont ceux des
problèmes des cas I, III, IV et VI dans la copie latine de ce volume
(vues 159, 175, 181, 187). La vue 302 photographie une seconde fois ce
même feuillet, cadré un peu autrement ; elle n'est pas transcrite. Le
verso est à la vue 303, blanc.}

\page{304}

On voit en effet que \struck{ses} \add{la valeur des} coëfficiens et l'équation qui détermine
la valeur d'$\omega$ dans le cas II. § 23 \add{d'une extremité libre et l'autre appuiée} sont d'accord avec les formules
relatives au cas I \struck{ou} des deux extremités libres \add{§ 15} car on a dans le
cas I lorsque $\text{sin}\,\omega$ est positif $\alpha = 1$, $\beta = -e^{\omega}$ $\gamma = 1+e^{\omega}$. $\delta = 1-e^{\omega}$
\uncertain{$A = -e^{\omega}$} \uncertain{$x = 1$}. et l'équation delaquelle on tire la valeur de l'angle
$\omega$ est $\text{cos}\,\omega = \frac{2}{e^{\omega}+e^{-\omega}}$. que j'ai déja montré etre equivallante a
$\text{tang}\,\frac{1}{2}\omega = \frac{1-e^{\omega}}{1+e^{\omega}}$ et pour le cas II $\alpha = 1$ $\beta = -e^{2\omega}$, $\gamma = 1+e^{2\omega}$, $\delta = 1-e^{2\omega}$
enfin l'équation qui détermine la valeur de l'angle $\omega$ est
$\text{tang}\,\omega = \frac{1-e^{2\omega}}{1+e^{2\omega}}$, il est clair que ces diverses quantités sont
identiques en observant \struck{que} de mettre 1 au lieu $\frac{1}{2}$ dans celles
qui concerne le cas II parceque la lame y est considérée comme
entière tandis que celle dont on détermine le mouvement d'après
la connoissance de celui dela lame dont les deux extremités
sont libres est censée etre la moitié de cette dernière.

On devroit trouver les memes rapports entre le cas V.
d'une extrémité fixée et l'autre appuiée § 39 et le cas VI.
des deux extremités fixées § 43. cependant on a dans le cas VI
\struck{\ill{}} lorsque $\text{sin}\,\omega$ est positif $\alpha = 1$, $\beta = -e^{\omega}$, $\gamma = -(1+e^{\omega})$, $\delta = -(1-e^{\omega})$
$\text{cos}\,\omega = \frac{2}{e^{\omega}+e^{-\omega}}$ ou $\text{tang}\,\frac{1}{2}\omega = \frac{1-e^{\omega}}{1+e^{\omega}}$, tandis que \struck{les valeurs du} \add{pour le}
cas V \struck{sont} \add{on a} \struck{\ill{}} alaverité $\text{tang}\,\omega = \frac{1-e^{2\omega}}{1+e^{2\omega}}$ mais que les coëfficiens
$\alpha = 1$, $\beta = -1$, $\gamma = \sqrt{2}\,(e^{2\omega}+e^{-2\omega})$ et $\delta = 0$ n'ont aucun rapport
apparans avec ceux du cas précédant, \struck{mais} et on a
\[ y = \ldots\, e^{x\omega} - e^{-x\omega} \mp \sqrt{2}\,(e^{2\omega}+e^{-2\omega})\,\text{sin}\,x\omega \]
tandis qu'il semble que
l'on devroit avoir
\[ y = \ldots\ldots\, e^{x\omega} - e^{(2-x)\omega} - (1+e^{2\omega})\,\text{sin}\,x\omega - (1-e^{2\omega})\,\text{cos}\,x\omega \]
Pour se rendre raison de ce paradoxe il faut remarquer
qu'Euler suppose que l'extrémité pour laquelle \ill{} $x = 0$ est
celle qui est appuiée, tandis que dans \uncertain{l'hypothese}
\note{En tête, au milieu, « 40 » souligné ; en haut à droite « 148 », à
l'encre. « la valeur » est écrit au-dessus de « ses », qui paraît biffé ;
« d'une extremité libre et l'autre appuiée » au-dessus de « cas II. § 23 »,
et « § 15 » en petit au-dessus de la fin de « libres ». Le groupe
\uncertain{$A = -e^{\omega}$} \uncertain{$x = 1$}, qui suit les quatre
coefficients du cas I, est tracé d'une encre très épaisse, comme repassé,
et sa lecture est douteuse lettre à lettre ; de même le groupe court,
biffé et empâté, qui ouvre la ligne après « cas VI », et celui qui suit
« sont » biffé après « cas V » (lecture impossible, un « \& » peut-être).
Dans « pour laquelle … $x = 0$ », un petit signe en forme de « t »
suivi de « ou » ou « on » précède $x = 0$ ; il n'est pas lu. Le dernier
mot, \uncertain{l'hypothese}, est traversé d'un trait et peut être
biffé ; la phrase s'arrête là, au bas du feuillet. Les références
§ 23, § 15, § 39 et § 43 sont celles des problèmes des cas II, I, V et VI
dans la copie latine de ce volume. Le verso (vue 305) est blanc.}

\page{306}

dans l'hypothese actuelle cene peut etre que pour une des
extremités dela lame entière et parconséquent pour l'extremité
fixée dela lame de longueur $\frac{1}{2}a$ que $x$ peut etre supposé
egal a zéro, il est clair au reste qu'il est fort indifferent
de faire l'une ou l'autre supposition.

En adoptant celle qui convient ici il faut mettre $1-x$ au
lieu de $x$ dans l'expression d'$y$ donnée par Euler elle
devient ainsi
\[ y = \ldots\ldots\, e^{(1-x)\omega} - e^{(x-1)\omega} \mp \sqrt{2}\{e^{2\omega}+e^{-2\omega}\}\,\text{sin}(1-x)\omega \]
\[ = \ldots\ldots\, e^{-\omega}\{e^{x\omega} - e^{(2-x)\omega} \pm \sqrt{2}\{e^{2\omega}+e^{-2\omega}\}\{\text{sin}\,\omega\,\text{cos}\,x\omega - \text{sin}\,x\omega\,\text{cos}\,\omega\}e^{\omega}\} \]
\uncertain{et} en mettant pour $\text{sin}\,\omega$ et $\text{cos}\,\omega$ leurs valeurs
$\pm\frac{(e^{\omega}-e^{-\omega})}{\sqrt{2}(e^{2\omega}+e^{-2\omega})}$, $\pm\frac{e^{\omega}+e^{-\omega}}{\sqrt{2}(e^{2\omega}+e^{-2\omega})}$ et observant que les signes superieurs
de ces valeurs et celui du terme $\pm e^{\omega}\sqrt{2}\{e^{2\omega}+e^{-2\omega}\}$ ont lieu
en même tems on a a cause de $e^{-\omega}$ constant
\[ y = \ldots\ldots\, e^{x\omega} - e^{(2-x)\omega} + (e^{2\omega}-1)\,\text{cos}\,x\omega - (e^{2\omega}+1)\,\text{sin}\,x\omega \]
et parconséquent $\alpha = 1$ $\beta = -e^{2\omega}$, $\gamma = -(1+e^{2\omega})$ $\delta = -(1-e^{2\omega})$
qui sont les valeurs que l'on auroit deduites de celles qui conviennent
au cas des deux extremités fixées.

Afin de faire sentir d'avantage combien cette déduction est fondée
on peut observer que les conditions indiquées par Euler N.o 13
pour que l'extremité d'une lame elastique vibrante soit simplement
appuiée sont $y = 0$ et $\left(\frac{ddy}{dx^2}\right) = 0$ \struck{ou $\left(\frac{ddy}{dx^2}\right) = 0$} et que
tant dans le cas des deux extremités libres que dans celui
des deux extremités fixées, on a lorsque $\text{sin}\,\omega$ est positif
$y = 0$ pour $x = \frac{1}{2}$ et que la valeur d'$y$ est alors exprimée
par $(1+e^{\omega})\,\text{sin}\,\frac{1}{2}\omega + (1-e^{\omega})\,\text{cos}\,\frac{1}{2}\omega = 0$ d'ou resulte evidemment
\[ \left(\frac{ddy}{dx^2}\right) = -(1+e^{\omega})\,\text{sin}\,\frac{1}{2}\omega - (1-e^{\omega})\,\text{cos}\,\frac{1}{2}\omega = 0. \]
\note{En tête, au milieu, « 41 » souligné, le chiffre repassé ; en haut à
droite « 149 », à l'encre. Le feuillet s'ouvre sur « dans l'hypothese »,
les deux mots par lesquels finit la vue 304 : elle les a répétés en tête
de la page suivante, et la phrase se lit sans rupture en passant de l'un
à l'autre. Dans la seconde ligne de $y$, la fin
« $\ldots\}e^{\omega}\}$ » est écrite serrée contre le bord et se lit
mal ; l'exposant de $e$ après « a cause de » est lu $-\omega$, et un
petit « $\omega$ » isolé au-dessus de « de » n'est pas rapporté. Le mot
qui ouvre la ligne « en mettant » est un signe en forme de croix, lu
\uncertain{et}. « N.o 13 » : le N est empâté d'une tache ; le § 13 de
la copie latine (lot 8) est la première scolie où Euler énonce les conditions
aux extrémités. La seconde condition biffée est barrée d'un trait
oblique. Le verso (vue 307) est blanc.}

\page{308}

Desorte que l'on est plainement autorisé a considérer lemilieu de
ces lames comme la jonction des extremités appuiées de deux autres
lames, il resulte dela que les ordonnées de ces dernières lames
devant rester les mêmes soit que l'on \struck{consi} les considere comme
nettement separées, ou que l'on les regarde comme fesant
partie dela lame primitive les coëfficiens qui entrent
dans l'expression de ces ordonnées sont également invariables
et qu'il suffit parconséquent de calculer ceux qui conviennent
au mouvement dela lame primitive pour avoir ceux
qui appartiennent a celui des lames dans lesquelles elle
peut etre censée divisée.

J'ai déja remarqué que la lame elastique dont les deux extremités
sont appuiées peut aussi etre censée partagée en deux ou plusieurs
autres lames elastiques de \struck{longueur} longueurs fractionnaires soumises
aux mêmes conditions, On seroit porté par analogie a croire que
dans le cas des deux extremités libres lorsqu'il n'y a pas de nœuds
au milieu les deux moitiés de la lame primitive peuvent
egalement etres considerées comme des lames entières soumises aux
mêmes conditions, mais les valeurs des ordonnées ne s'accomodent point
de cette hypothèse et on en voit la raison dans l'observation
que j'ai faits delaquelle il resulte que lamoitié dela
partie comprise entre les deux nœuds les plus voisins du milieu de
la lame primitive ne peut etre egale aux parties comprises entre
les nœuds les plus voisins des extremités de cette lame et les mêmes
extremités
\note{En tête, au milieu, « 42 » souligné ; en haut à droite « 150 », à
l'encre. Après « que j'ai faits » la ligne laisse un blanc d'un tiers
environ avant « delaquelle », comme pour un renvoi à remplir ; il n'est
pas rendu ici. Le texte s'arrête sur « extremités », sans point ;
le bas du feuillet est blanc. Le verso
(vue 309) est blanc.}

\page{310}

Les vibrations qui font l'objet du mémoire d'Euler et
des \struck{reflexions} remarques précedantes sont celles que M\textsuperscript{r} Chladni
nomme vibrations transversales dela lame; Mais les vibrations des
plaques elastiques \struck{dont les} libres dans le cas ou il y a des lignes
nodales en \struck{differens} suivant differentes directions, dont la théorie
semble devoir etre beaucoup plus compliquée presente cependant
une singularité remarquable c'est que tandis que pour les
vibrations transversales dela lame elastique dont les extremités
sont libres, on ne peut composer une lame de deux moitiés
soumises aux mêmes conditions, il y a au contraire plusieurs
cas ou les \struck{lames} plaques carrées vibrantes en differens sens a la
fois, et libres de tous cotés, peuvent cependant etre formées
de quatre autres lames soumises aux mêmes conditions voyez
traité d'Acoustique par Chladni §§ 110, en comparant les
differentes figures qui presentent entr'elles le rapport dont
il s'agit on voit que les figures composées sont celles
ou il n'y a pas de lignes nodales qui coupent \struck{la surface
en deux parties egales de mi} diamétralement la surface
en deux parties égales, et on sent aisement que cela doit
etre en effet puisque pour qu'une partie continue puisse
equivaloir a une extremité ou coté libre il faut qu'elle
soit egalement libre \struck{c'est adire en mouvement}; et parconséquent
en mouvement: par exemple, la fig. 63 qui est la plus simple
parmi celles qui presentent des lignes nodales en differens sens,
repetée quatre fois forme la fig 64. a et l'analogie \struck{porte}
\struck{porte a croire} \add{montre} que la même figure 63 repetée 9 fois donneroit
la figure 75 et \add{que cette même fig 63} prise \struck{huit} \add{16 fois, formeroit} la fig. 82, \add{\struck{quoique}} \struck{mais cependant M\textsuperscript{r}
Chladni ne dit pas avoir obtenu ces deux dernieres figures par
le rapprochement de 9 et 16 plaques portant la fig. 63}
\note{En tête, au milieu, « 43 » souligné ; en haut à droite « 151 », à
l'encre. Le renvoi à Chladni est lu « §§ 110 ». Dans « la fig. 63 qui
est la plus simple », une tache couvre la fin de « fig » ; le chiffre 63
est net. « 64. a » est écrit ainsi. Les additions « montre », « que cette
même fig 63 » et « 16 fois, formeroit » sont au-dessus de la ligne, sans
signe d'insertion ; « huit » est biffé d'un trait épais. Au-dessus de
« mais » biffé, « quoique » est écrit puis biffé à son tour. Les deux
dernières lignes et la fin de la précédente, depuis « mais cependant »,
sont barrées ligne à ligne ; la phrase reste ainsi sans fin au bas du
feuillet. Le verso est à la vue 311.}

\page{312}

\struck{Mais si cette conjecture etoit sanctionnée par l'expérience on seroit}
\struck{Il me semble que l'on peut conclure de cette analogie qu'il y a aumoins
quelques cas ou la fonction qui en supposant la plaque carrée
on peut supposer, que les coordonnées dela surface vibrante sont symétriques
autour des}

\struck{Il me semble que l'on peut en inferer que} \add{Ainsi} dans certains cas au moins
si on suppose une \struck{surface} \add{plaque} vibrante elastique quarrée dont la superficie
soit exprimée par $aamm$ et que l'on \struck{applique un stilet} divise ce
quarré en $mm$ quarrés degrandeurs $aa$ et que l'on applique un stilet
au point convenable et sous les conditions propres a faire naitre dans
un des 4, quarrés \add{placés aux angles dela plaque} d'etendue $aa$ une certaine figure nodale cette même
figure sera repetée dans les \uncertain{$mm-1$} autres carrés qui composent la plaque
entière \struck{desorte que dans l'exemple dont il s' c'est adire} \add{d'ou il resulte que le mouvement \struck{parconséquent} de ces quarrés elementaires est \struck{uniforme} semblable dans chacun d'eux et parconséquent} que si on
suppose \struck{un} \add{deux} plans perpendiculaires \add{entr'eux et l'un et l'autre perpendiculaires au plan} \struck{a} l'horizontal qui est \struck{enmême tems} \add{celui}
\add{de} la figure \struck{initiale} \add{\struck{no}} dela plaque \add{en repos}, les courbes formées par \add{les} \struck{l'}intersection
de ces plans, \add{supposés encore perpendiculaires aux cotés de cette même plaque quarrée} et de la surface vibrante couperont le plan \struck{horiz}
horizontal dans des points dont les \struck{abscisses seront} coordonnées
seront $\alpha$, $\alpha+a$, $\alpha+2a$ \struck{\ill{}} et $\alpha'$, $\alpha'+a$, $\alpha'+2a$ \uncertain{\&c.}
\struck{\ill{}} \add{$\alpha$ et $\alpha'$ etant} dependant dela nature dela figure
\struck{elemen} formée dans le \struck{plus} quarré elementaire, et que \struck{ces}
courbes \struck{seront} symétriques autour des points dont il s'agit.

Cette propriété \struck{que M\textsuperscript{r} Chladni etablit par rapport a} pour
\add{des plaques quarrées elastiques vibrantes libres de tous cotés de pouvoir etre considerées comme}
\struck{differentes figures nodales dans le cas ou m=2} d'etre exactement décom-
\add{decomposées en un certain nombre de plaques de même forme soumises \struck{aux} mêmes}
posable en quatre \struck{figures} plaques chargées de figures \struck{semblables}
\add{conditions, propriétés que M\textsuperscript{r} Chladni}
\marginal{établit dans plusieurs cas pour $m=2$, \struck{et} que je viens demontrer \struck{\ill{}} \add{avoir} encore lieu \struck{pour} \add{dans} le cas de la fig 63. pour $m=3$ et $m=4$}
\struck{pour M\textsuperscript{r} Chladni etablit pour differentes figures nodales dans}
\add{\struck{et que}} \struck{le cas ou $m=2$}, conduit donc a \struck{fin} conclure qu'il
existe dans certains cas degrandes analogies entre les vibrations
des plaques elastiques quarrées et celle des surfaces tendues demême
figure. \add{car} on peut voir \add{par la comparaison du § \uncertain{26} du mémoire de M\textsuperscript{r} Biot} \struck{dans le mémoire de M\textsuperscript{r} Biot \ill{}}
inseré dans le tome 4 des memoires de l'institut mémoire
dans lequel l'auteur \struck{détermine} donne l'équation des surfaces
\add{il n'est question que de ce genre de surface}
vibrantes \struck{tendues}, que dans ce cas $\alpha$ et $\alpha'$ sont \struck{zero} \add{zero ce qui doit etre en effet} parceque
\add{dans \struck{aucune} ces surfaces nécessairement}
les lignes de repos sont placées aux extremités des surfaces composantes
\marginal{qui peuvent etres considerées comme vibrant separement.}
\marginal{on voit encore que pour les figures composées dela fig 63 $\alpha$ et $\alpha' = \frac{1}{2}a$}
\note{En tête, au milieu, « 44 » souligné ; en haut à droite « 152 », à
l'encre. Brouillon très corrigé. La première ligne est barrée seule ; les
quatre suivantes (« Il me semble… autour des ») sont traversées chacune
d'un long trait qui court d'un bout à l'autre et sont données comme
biffées. Les additions sont écrites entre les lignes d'une plume plus
fine, sans signe d'insertion ; elles sont placées ici avant les mots
qu'elles surmontent, et l'ordre de lecture de certaines, dans les lignes
« d'ou il resulte… parconséquent » et « Cette propriété… conduit donc », est
celui que suggère le sens, non un ordre marqué sur la page. Dans ce
dernier passage, deux rédactions se superposent : la première,
« Cette propriété que M\textsuperscript{r} Chladni etablit par rapport a
differentes figures nodales dans le cas ou m=2 d'etre exactement
décomposable en quatre plaques chargées de figures… », est biffée en
partie seulement ; la seconde s'écrit dans l'interligne et se poursuit
dans la marge de gauche (« établit dans plusieurs cas pour $m=2$… pour
$m=3$ et $m=4$ »), en face des lignes qu'elle complète. De même la note
de marge « qui peuvent etres considerées comme vibrant separement. »
achève la dernière ligne du feuillet (« des surfaces composantes ») ; la
ligne qui la suit dans la marge, d'une plume plus fine, est une
remarque à part. Dans la marge de gauche encore, en face de « formée dans
le quarré elementaire », un petit dessin : un carré partagé en
quatre par deux lignes médianes, traversé d'une diagonale, avec un petit
arc dans le carré du haut à gauche. La
ligne des coordonnées porte après $\alpha+2a$ un groupe biffé, et la
ligne suivante s'ouvre sur une suite de termes biffés d'un trait épais,
non lus ; « \&c. » après $\alpha'+2a$ est lu avec doute (« de. »).
« $mm-1$ » : un trait et une virgule suivent « mm », lus comme $-1$ avec
doute. « § 26 » : le signe § est empâté et le 6 a la forme d'un « b »,
comme ailleurs dans ce volume. La mention de Biot et du « tome 4 des
memoires de l'institut » est de sa main ; le passage n'a pas été vérifié
contre ce mémoire. Dans la dernière note de marge, le dénominateur de
$\frac{1}{2}a$ est petit et serré et se lit mal. Le verso est à la vue 313.}

\page{314}

La théorie complette dela plaque elastique vibrante quarrée \struck{si on avoit
egard a l'etat des cotés qui pourroient etre appuiés} comprendroit
quinze cas differens, si on avoit égard a l'etat des cotés qui pourroient
etre appuiés libres ou fixés, mais il arriveroit sans doute comme pour les
vibrations transversales delalame elastique que ces cas auroient entr'eux
des relations qui permettroient de n'en calculler que quelques uns, par exemple
le cas ou une figure reguliere \struck{de trois} formée par les lignes nodales dans
une plaque quarrée dont tous les cotés sont libres presente deux ligne
diametrales qui se coupent aucentre de cette figure pourroit peut etre
etre considéré comme produit par la reunion dequatre plaques
quarrées degrandeurs convenables qui auroient chacune deux deleurs
cotés libres et les deux autres appuiés.

\struck{\ill{}} Entre les figures regulieres l'ouvrage de M\textsuperscript{r} Chladni
en presente un grand nombre dans lesquelles il n'y a aucun
rapport apparent entre une première moitié et la seconde.
je crois que l'on peut penser que dans ce dernier genre de figures
le point d'appui qui appartient a une des ligne nodale
gêne plus ou moins le mouvement qu'il ne l'est dans les autres lignes
derepos; il me semble en effet que \struck{ce n'est pas assez des effets}
si le mouvement étoit egalement libre dans toute l'étendue dela
plaque il ne devroit se former que des figures régulières.

\struck{Que ce soit que dans l'espece dont il s'agit ici il puisse exister}

Cette supposition \struck{qui} me semble devoir convenir ici parceque lorsmême qu'on
admettroit que le mouvement seroit plus gêné dans le point d'appui que
dans les autres points qui enforment les lignes nodales \struck{sont censés}
sont censés equivallant \struck{a} aux points d'appui, la communication entre
les differentes parties dela plaque vibrante ne seroit pas pour cela
interrompue puisqu'independamment de ce point d'appui pris \struck{les} parties
vibrantes endeça et audela, \struck{communiquent} se touchent par les cotés
qui leurs servent delimites. \struck{On voit assez que dans lorsqu'il est question
des vibrations transversales dela lame, on ne peut supposer que le
point d'appui interrompe la communication du mouvement soit interrompue
\ill{} le point d'appui sans admettre qu'elle soit entièrement detruite}
\marginal{\struck{que je crois que pour ce genre} \struck{car on fait toujours abstraction dela largeur de ces lames, ainsi pour ce genre de vibration je crois que l'on peut toujours regarder le point d'appui comme equivallant}}
\note{En tête, au milieu, « 45 » souligné ; en haut à droite « 153 », à
l'encre. À la première ligne, « plaque », « elastique », « vibrante » et
« quarrée » portent chacun à la fin une petite croix, qui paraît biffer un
« s » du pluriel ; le singulier est donné. Le mot biffé qui ouvre le
second paragraphe, devant « Entre », n'est pas lu. La ligne « Que ce soit
que… exister » est barrée d'une suite de hachures obliques. Après
« delimites. » un signe « + » renvoie à la note de la marge de gauche ;
cette note (sa première ligne barrée d'un trait, le reste de hachures
obliques) et les quatre dernières lignes du texte, depuis « On voit
assez », sont annulées ensemble par une même série de hachures ; le mot
qui ouvre la dernière ligne n'est pas lu. Le bas du feuillet est abîmé
et replié. Le verso est à la vue 315.}

\page{316}

si on vouloit avoir la théorie complette des plaques quarrées
elastiques vibrantes en ayant égard a l'état des cotés qui
peuvent etres appuiés fixés ou libres on auroit 15 cas
differens suivans qu'un ou plusieurs de ces cotés seroient
dans l'un ou l'autre des trois etats \struck{dont que}
l'on distingue mais il arriveroit sans doute comme
dans la théorie des lames elastiques vibrantes
que ces \struck{etats} \add{cas} auroient entr'eux des relations
qui permettroient de dans \struck{\ill{}} calculler que quelquesuns
Cette immensité de cas dont ceux que M\textsuperscript{r} Chladni
afait connoître ne sont qu'une petite partie
se trouve encore augmenté \struck{par la remarque}
d'une manière effrayante lorsque l'on remarque
qu'il parle dans plusieurs experiences de degrés
differens de pression necessaire pour la formation
de telle ou telle figure car on voit alors
que les etats distingués par euler n'ont pas
dans la fixité que \struck{\ill{}} qu'il leur suppose
dans la théorie car une pression foible differe
peu de l'état de liberté \struck{et on peut la}
supposer croître par \struck{une} une échelle composée
d'un nombre infini de degré dont la limite
est la fixité\textsuperscript{+} \struck{le nombre de variations possibles
seroit encore augmenté si on vouloit tenir compte
des position et pression et vivacité du mouvement
de l'archet que circonstances qui contribuent tous
a differencier les experiences}
\marginal{+ on concoit en effet que dans la vibration des plaque la fixité du stilet \struck{ne s'op} doit plus empecher la communication du mouvement \struck{ex} entre les portions \struck{de la} placées de part et d'autre de cet obstacle puisqu'il y a continuité dans les points qui restent libres pendant son action}
\note{Pas de numéro en tête ; en haut à droite, décalé vers la gauche,
« 154 », à l'encre. Ce feuillet reprend, dans une autre rédaction, le
début de la vue 314 (« La théorie complette dela plaque elastique
vibrante quarrée… quinze cas differens ») et le poursuit autrement ;
rien sur le feuillet ne dit laquelle des deux est la première. Le texte
est écrit en lignes plus courtes, la marge de gauche portant une note
longue, appelée par le « + » qui suit « la fixité » ; elle est donnée
ici après le texte. « de dans » est écrit ainsi, là où le sens demande
« de n'en » ; après « dans », un signe court biffé n'est pas lu, non plus
que le mot court biffé après « la fixité que ». « euler » est écrit sans
majuscule. Les sept dernières lignes, depuis « le nombre de variations »,
sont barrées ligne à ligne. Le verso (vue 317) est blanc.}

\page{318}

\begin{quote}
Memoire presenté et \uncertain{examiné}

sur la lame
\end{quote}
\note{Feuillet de titre, qui ne porte que ces deux lignes, au tiers
supérieur de la page, la seconde soulignée d'un long trait ; en haut à
droite « 155 », à l'encre. Le papier est d'un ton plus sombre que celui
des feuillets qui précèdent. La fin de « examiné » est écrite sur d'autres
lettres, repassée d'une encre plus épaisse ; la lecture est offerte avec
doute. Rien d'autre sur le feuillet ; le verso (vue 319) est blanc.}

\page{320}

\section{Remarques sur le mémoire d'Euler: Investigatio motuum quibus laminæ e virgæ elasticæ contremiscunt. in act acad Petrop ann. 1779 P 1 p. 103 et seq}

Après avoir examiné les mouvemens reguliers possibles dela lame elastique vibrante,
en ayant seulement égard à l'état des extremités qui peuvent etre fixées, simplement
appuyées ou libres, cequi donne lieu à six cas essentiellement differens, l'auteur suppose
que la même lame puisse en outre etre appuyée dans un ou plusieurs points differens
des extremités, de façon \struck{cequ} que la communication entre ces differentes parties nesoit pas
interrompue; et comme le developpement de tous les cas demanderoit des recherches
infinis, il se contente d'examiner un d'entr'eux, dans l'intention de faire comprendre
comment le calcul peut leur etre appliqué. celui qu'il choisit qui est le plus
simple de tous, est le cas où la lame soumise à l'action du stilet, est appuyée
a l'une et l'autre de ses extremités. il se propose, § 47, de trouver tous les
mouvemens de vibration dont une telle lame peut etre agitée.

Dans la suite du même §, l'auteur parvient à l'équation suivante
\[ 0 = 2 - 2e^{2\omega} - (e^{x\omega} - e^{\omega(2-x)})(e^{x\omega} - e^{-x\omega})\,\frac{\text{sin}\,\omega}{\text{sin}\,x\omega\;\text{sin}(1-x)\omega}. \]
il remarque que la question est réduite
à tirer de cette équation les valeurs de l'angle $\omega$: mais il ajoute qu'il n'ose
entreprendre cetravail à cause de la complication des formules, et qu'il suffit
d'avoir montré dans celieu par qu'elle méthode on peut traiter ces difficiles
questions. Cependant, après avoir établi que les suppositions $x = 0$ et $x = 1$ menent à
conclure $\omega = 0$, qui est la valeur qui convient au repos dela lame, il observe que
le cas où $x = \frac{1}{2}$ merite d'être développé. L'équation finale se partage \add{dans ce cas} endeux
facteurs qui, egalés séparement a zéro, donnent lieu à deux solutions analisées
§ 49 et 50 pour la première, § 51 pour la seconde; et l'auteur termine en remarquant
qu'il est digne d'attention que les deux moitiés dela lame puissent vibrer dedeux
manières, l'une qui se rapporte au cas IV, et l'autre qui appartient au cas V.
\note{En tête, au milieu, « (1) » souligné ; en haut à droite « 156 », à
l'encre. Une pièce nouvelle commence ici, en belle copie, d'une écriture
régulière et presque sans rature : le titre sur deux lignes, souligné d'un court trait au milieu, puis un
exposé à la troisième personne (« l'auteur ») du problème du § 47
d'Euler, la lame appuyée aux deux bouts et soumise au stilet. « e virgæ »
est écrit ainsi pour « et virgæ » ; « P 1 » pour la première partie du
volume. « p. 103 » est la page par laquelle s'ouvre aussi la copie latine
de ce volume (vue 144). L'équation est écrite dans la ligne, suivie d'un
point et de « il remarque… » ; elle est détachée ici. Les renvois
« § 49 et 50 » et « § 51 » sont lus sur la page, le 4 tracé comme ailleurs
à la façon d'un « b ». Le texte occupe les deux tiers du feuillet et
s'arrête sur « au cas V. », le bas étant blanc. Un cachet de la Bibliothèque royale se devine, très pâle, à droite de la
page : c'est celui du verso, vu par transparence. Le verso (vue 321,
hors de ce lot) est blanc à la vignette, hors ce cachet.}

\end{document}
